\documentclass[11pt]{article}

% ---------------------------------------------------------------------------
% Learning to Deceive and Persuade: Predictive Agents with Language and
% Personality Diversity in Social Deduction Games.
%
% Self-contained draft (compiles with a stock TeX install; no venue .sty
% required). Swap \documentclass for the target venue style before submission.
% ---------------------------------------------------------------------------

\usepackage{arxiv}          % self-contained arXiv-preprint style (see arxiv.sty)
\usepackage{amsmath,amssymb}
\usepackage{booktabs}
\usepackage{multirow}
\usepackage{graphicx}
\usepackage{xcolor}
\usepackage[numbers,sort&compress]{natbib}
\usepackage[hidelinks]{hyperref}

% Draft note macro — remove before submission.
\newcommand{\note}[1]{\textcolor{red}{[#1]}}
\newcommand{\ci}[2]{{\footnotesize [#1, #2]}}
% \num as a no-op so the draft compiles without siunitx; swap in siunitx later if desired.
\newcommand{\num}[1]{#1}
% \TBD marks a numeric cell to be filled from the pending publish-run outputs.
\newcommand{\TBD}{\textcolor{red}{--}}

\title{Learning to Deceive and Persuade: Predictive Agents with\\
Language and Personality Diversity in Social Deduction Games}
\author{Aaron Shah\\ Virginia Commonwealth University}
\date{\today}

\begin{document}
\maketitle

\begin{abstract}
We study whether a joint-embedding predictive architecture (JEPA) can learn the
strategic core of a social-deduction game---one-day \textit{Werewolf} with nine
players---and how much natural language contributes on top of a learned latent
policy. Each agent maintains a belief latent produced by a shared encoder, a
world model that predicts the next latent, and a factorized planner with
separate talk, vote, and kill heads. A social-influence module applies a
wolf-supervised correction to the belief latent, and a language ``mouthpiece''
biases a frozen language-model decoder toward a chosen intent and target while an
LLM-as-judge gates message acceptance. We evaluate the system with a
\textit{retrain-per-condition} ablation ladder: every rung is scored against a
model trained with exactly the components it is evaluated with, so no ablation is
judged on a policy that never saw its own inputs. Our central finding is that
language and the learned policy are \emph{not separable wins}: the full system
nearly doubles the villager win rate over the strongest non-language ablation
(\num{0.485} vs.\ \num{0.303}) and this language$+$judge contrast is the only
statistically significant rung of the ladder, yet language paired with
\emph{random} voting performs no better than random play. We further show that
training the social-influence correction jointly with the planner reverses an
apparent ``social hurts'' effect observed under a train/evaluate mismatch, and
yields a markedly sharper world model.
\end{abstract}

% ===========================================================================
\section{Introduction}
% ===========================================================================
Social-deduction games such as \textit{Werewolf} (\textit{Mafia}) demand exactly
the capabilities that are hardest to isolate in language agents: maintaining a
belief over hidden roles, reasoning about how one's own utterances shift others'
beliefs, and converting that reasoning into a coordinated vote. Recent work
drives such agents with large language models end to end, which entangles
strategy and language and makes it difficult to ask how much of the competence is
\emph{strategic} versus \emph{rhetorical}~\citep{xu2023werewolf}.

We take the opposite decomposition. Following the predictive-architecture program
of \citet{lecun2022path} and \citet{assran2023ijepa}, we place a small,
\emph{learned} latent policy at the center of the agent---a belief encoder, a
world model, and a factorized planner---and treat language as a controllable
channel layered on top rather than the seat of strategy. This lets us pose the
question the entangled designs cannot: \textbf{does language help because it
conveys strategy the policy already computes, or is language a free win on its
own?}

Answering it cleanly requires an ablation methodology that most such studies skip.
If the policy is trained once with all components and then evaluated with
components switched off, an ablation measures a \emph{distribution shift} away
from training, not the causal contribution of the component. We therefore adopt a
\textbf{retrain-per-condition} ladder: each rung is trained with exactly the
components it will be evaluated with. This single methodological change reverses
one of our own earlier conclusions (Section~\ref{sec:results}).

\paragraph{Contributions.}
\begin{enumerate}
  \item A JEPA-centered Werewolf agent with a factorized talk/vote/kill planner, a
        jointly-trained social-influence correction, and a language mouthpiece
        gated by an LLM judge (Section~\ref{sec:method}).
  \item A retrain-per-condition ablation ladder (B0--B6) that scores every rung
        against a model trained with its own components
        (Section~\ref{sec:setup}).
  \item Evidence that language and the learned policy are inseparable: the full
        system nearly doubles the villager win rate over the best non-language
        ablation and is the only significant rung, while language with random
        voting equals random play (Section~\ref{sec:results}).
\end{enumerate}

% ===========================================================================
\section{Related Work}
% ===========================================================================
Our design sits at the intersection of four lines of work: predictive
architectures and world models, multi-agent learning in games of hidden
information, language and deception in interactive agents, and machine theory of
mind. We also position our ablation methodology against the common practice of
train-once/evaluate-many.

\paragraph{Predictive architectures and world models.} JEPA-style models learn by
predicting a target \emph{representation} of future input in a latent space rather
than reconstructing raw observations~\citep{lecun2022path,assran2023ijepa}. This
lineage is closely related to model-based reinforcement learning that plans in a
learned latent---World Models~\citep{ha2018world} and latent-imagination agents
such as Dreamer~\citep{hafner2020dream}---and to the free-energy view of
perception and action as prediction-error
minimization~\citep{friston2010free}. Where prior work uses the learned dynamics
model chiefly for control, we additionally repurpose one-step latent prediction
error as an intrinsic \emph{value} signal: our language-free B3 baseline votes for
the target whose elimination the world model predicts will most reduce future
prediction error, giving a policy-head-free lower bound on what the JEPA latent
alone supports. We use it as a language-free voting rule.

\paragraph{Multi-agent learning in games of hidden information.} Games with
partial observability and hidden roles have become standard testbeds for
coordination and inference. Hanabi stresses implicit signalling under a no-talk
constraint~\citep{bard2020hanabi}, and the Bayesian Action Decoder learns
conventions by reasoning over public
beliefs~\citep{foerster2019bayesian}. Closest to Werewolf, DeepRole combines
counterfactual regret minimization with explicit deductive reasoning over hidden
roles in \emph{The Resistance: Avalon}~\citep{serrino2019finding}. These systems
achieve strong play with hand-designed or tightly-structured inference; our
contribution is orthogonal---a \emph{learned} latent belief and factorized planner
whose components we can switch on and off---and we ask how much a natural-language
channel adds on top of that learned core.

\paragraph{Language, negotiation, and deception.} Emergent-communication work shows
that agents can learn discrete protocols to coordinate under partial
observability~\citep{foerster2016learning}, and end-to-end negotiation agents learn
to bargain---and to bluff---directly from dialogue~\citep{lewis2017deal}. At larger
scale, CICERO couples a strategic planner with a language model to reach
human-level play in Diplomacy, grounding messages in intended
actions~\citep{fair2022diplomacy}. We adopt the same planner/language separation
but invert the emphasis: rather than optimizing dialogue for persuasion, we hold
the language channel fixed (an intent-conditioned mouthpiece over a frozen decoder,
gated by a judge) and measure its \emph{marginal} contribution to win rate against
a learned policy. This lets us test whether fluent, judge-approved speech is worth
anything when the underlying voting policy is uninformed---a question the
end-to-end systems cannot isolate.

\paragraph{LLM social-deduction agents.} A recent line drives Werewolf and Avalon
agents with large language models, either prompting frozen models to
reason~\citep{xu2023werewolf} or adding adaptive/experience-based mechanisms to
improve strategic play~\citep{lan2024werewolf}. More broadly, generative agents use
LLMs to produce believable social behavior in open-ended
simulations~\citep{park2023generative}. Because these agents entangle strategy and
language inside one model, strong win rates are difficult to attribute: our
decomposition and the B4 (language$+$random-voting) rung are designed precisely to
separate rhetorical fluency from strategic competence
(Section~\ref{sec:results}).

\paragraph{Theory of mind and belief modeling.} Modeling other agents' hidden
mental states is central to social deduction. Machine-theory-of-mind networks learn
to predict others' beliefs and behavior from observation~\citep{rabinowitz2018machine},
and Bayesian theory-of-mind underlies DeepRole's role
inference~\citep{serrino2019finding}. Our social-influence module is a lightweight,
learned analogue: it applies a wolf-supervised correction to an agent's belief
latent conditioned on received messages, and---crucially---is trained
\emph{jointly} with the planner so the policy learns to use it rather than being
perturbed by it at test time.

\paragraph{Representation collapse and LLM-as-judge.} Latent-predictive training can
collapse to constant representations; variance/covariance
regularization~\citep{bardes2022vicreg} and normalization~\citep{ba2016layernorm}
counteract this. We found layer normalization on the encoder output essential for
stable latents, after a variance-hinge regularizer alone failed
(Section~\ref{sec:method}). For message quality we use an
LLM-as-judge~\citep{zhao2023llmjudge} that scores coherence, truthfulness,
role-consistency, and safety, and note the standard caveat that a judge sharing a
model family with the speaker tends toward leniency (Section~
ef{sec:discussion}).

\paragraph{Ablation under train--evaluate consistency.} Finally, our study is a
methodological cautionary tale. Component ablations are routinely computed by
training a single full model and disabling parts at evaluation. When a component
(here, the social correction) participates in the learned policy, disabling it at
test time measures a distribution shift away from training rather than the
component's causal value---and, as we show, can invert the \emph{sign} of the
measured effect. We therefore retrain per condition, an inexpensive discipline that
we argue should be standard for component-level claims about interactive agents.

% ===========================================================================
\section{Method}
\label{sec:method}
% ===========================================================================
\paragraph{Game.} We play a nine-player one-day-cycle Werewolf variant: seven
villagers and two werewolves. Each round has a discussion phase (agents emit
messages), a day vote (strict majority; a tie yields no elimination), and a night
kill by the werewolf team (target $i$ drawn with probability
$p_i \propto \exp(c_i / \max_j c_j)$ over suspicion scores $c$). The game ends when
all werewolves are eliminated (villagers win) or werewolves reach parity.

\paragraph{Belief encoder.} A shared multilayer-perceptron encoder maps each
agent's observation to a belief latent $z \in \mathbb{R}^{32}$. The observation
includes public game state, per-player features, and an \emph{identity block}---a
role bit, a self one-hot, and a persona embedding---so that the latent is
conditioned on \emph{who} is observing. We apply layer
normalization~\citep{ba2016layernorm} to the encoder output, which bounds
$\lVert z \rVert$ and was necessary to prevent both representational collapse
(cosine similarity $\to 1$) and norm explosion during joint training.

\paragraph{World model and factorized planner.} A world model $g$ predicts the
next belief latent $\hat z_{t+1} = g(z_t, a_t)$; targets are produced by
re-encoding the realized next observation with a stop-gradient, giving a JEPA
prediction loss $\lVert \hat z_{t+1} - \text{sg}(z_{t+1}) \rVert^2$. A
\emph{factorized planner} maps the belief latent to three action heads---talk
(intent/target), vote, and kill---trained by behavior cloning on rollouts. A
language-free baseline (B3) votes by free energy: it scores candidate targets by
the one-step prediction error they induce, using no planner head.

\paragraph{Social-influence correction.} A social module produces a correction
$\delta$ to the belief latent, relative-scaled so
$\lVert \delta \rVert = s \cdot \lVert z \rVert$, and conditioned on the messages
an agent has received. We consider two training regimes. In the \emph{detached}
regime, $\delta$ is supervised to point toward werewolves and is not backpropagated
into the planner---this is the natural but, as we show, \emph{harmful} default. In
the \emph{integrated} regime the planner's behavior cloning is trained on
$z + \delta$ (with $\delta$ in the gradient path), so the policy learns to
\emph{use} the social signal. Unless noted, ``social'' refers to the integrated
regime.

\paragraph{Language mouthpiece and judge.} Rather than let a language model choose
actions, we decouple content from strategy. The planner emits a discrete intent
(accuse, question, defend, vote, hedge) and a target; a \emph{mouthpiece} biases a
frozen decoder's logits toward realizing that intent, producing a natural-language
message. An LLM-as-judge~\citep{zhao2023llmjudge} scores each message on
coherence, truthfulness, role-consistency, and safety, and gates acceptance. This
keeps strategy in the learned policy while language remains fluent and auditable.

% ===========================================================================
\section{Experimental Setup}
\label{sec:setup}
% ===========================================================================
\paragraph{Retrain-per-condition ladder.} We evaluate seven rungs, from the full
system (B0) down to random play (B6). Crucially, each rung is scored against a
model \emph{trained} with exactly the components it is \emph{evaluated} with. The
seven rungs are covered by four deduplicated training configurations (a
no-social/no-language base shared by B2/B3/B5/B6; an integrated-social model for
B1; a language model with random voting for B4; and the full model for B0), each
written to its own checkpoint namespace. Table~\ref{tab:rungs} lists the rungs.

\begin{table}[t]
\centering
\small
\begin{tabular}{lll}
\toprule
Rung & Components (trained \& evaluated) & Voting policy \\
\midrule
B0\_full            & JEPA + planner + social + language + judge & planner \\
B1\_planner\_social & JEPA + planner + integrated social         & planner \\
B2\_jepa\_planner   & JEPA + planner                             & planner \\
B3\_jepa\_only      & JEPA world model (no planner head)         & free-energy \\
B4\_llm\_only       & language + judge                           & random \\
B5\_heuristic       & none                                       & heuristic \\
B6\_random          & none                                       & random \\
\bottomrule
\end{tabular}
\caption{The ablation ladder. Each rung is trained and evaluated with the same
component set (retrain-per-condition).}
\label{tab:rungs}
\end{table}

\paragraph{Metrics.} We report villager win rate (primary), villager vote accuracy
(fraction of villager votes cast at a werewolf), one-step world-model prediction
error ($\Delta z$ MSE), and judge acceptance. We compute 95\%
bootstrap confidence intervals over per-game outcomes and report contrasts as
differences in win rate with a bootstrap CI and $P(\Delta>0)$.

\paragraph{Scales.} We report a \emph{medium} run (99 games per rung across three
seeds) that resolves the primary contrast, and note where a full-scale run
(\num{450} games $\times$ 3 seeds, plus persona/social sweeps) is pending
(Section~\ref{sec:results}). Language rungs (B0, B4) use \texttt{gpt-4o-mini} as
both speaker and judge.

% ===========================================================================
\section{Results}
\label{sec:results}
% ===========================================================================
Table~\ref{tab:table1} reports the full ladder at medium scale. Three findings
stand out.

\begin{table}[t]
\centering
\small
\begin{tabular}{lcccc}
\toprule
Rung & Win rate \ci{95\%}{CI} & Vote acc. & $\Delta z$ MSE & Judge acc. \\
\midrule
\textbf{B0\_full}   & \textbf{0.485} \ci{0.384}{0.586} & 0.301 & 0.594 & 0.998 \\
B3\_jepa\_only      & 0.303 \ci{0.212}{0.394}          & 0.301 & 4.08  & --    \\
B1\_planner\_social & 0.222 \ci{0.141}{0.303}          & 0.249 & 0.583 & --    \\
B2\_jepa\_planner   & 0.202 \ci{0.121}{0.283}          & 0.307 & 3.79  & --    \\
B4\_llm\_only       & 0.192 \ci{0.121}{0.273}          & 0.294 & 4.78  & 0.999 \\
B6\_random          & 0.192 \ci{0.121}{0.273}          & 0.294 & 3.89  & --    \\
B5\_heuristic       & 0.071 \ci{0.030}{0.131}          & 0.289 & 3.81  & --    \\
\bottomrule
\end{tabular}
\caption{Table 1 --- ablation ladder, medium scale ($n=99$ per rung, 3 seeds).
Win rate is the villager win rate. \note{Replace with full-scale numbers
($n=450\times3$) from the pending run; add Table 2 (persona/social sweeps).}}
\label{tab:table1}
\end{table}

\begin{table}[t]
\centering
\small
\begin{tabular}{lcc}
\toprule
Contrast & $\Delta$ win rate \ci{95\%}{CI} & Verdict \\
\midrule
Language$+$judge \; (B0 $-$ B1) & $+0.263$ \ci{+0.131}{+0.394} & \textbf{significant} ($P{=}1.00$) \\
JEPA$+$planner \; (B2 $-$ B6)   & $+0.010$ \ci{-0.101}{+0.121} & n.s. \\
Social \; (B1 $-$ B2)           & $+0.020$ \ci{-0.091}{+0.131} & n.s. \\
\bottomrule
\end{tabular}
\caption{Ladder contrasts (medium scale). The language$+$judge contrast is the
only significant rung.}
\label{tab:contrasts}
\end{table}

\paragraph{(1) Language and the learned policy are inseparable.} The full system
(B0) wins \num{0.485} of games---nearly double the strongest non-language rung
(B3, \num{0.303}) and more than double random (\num{0.192}). The language$+$judge
contrast $\mathrm{B0}-\mathrm{B1} = +0.263$ is the only significant rung of the
ladder ($P(\Delta>0)=1.00$; Table~\ref{tab:contrasts}). Yet language is not a free
win: \textbf{B4 (fluent language with a \num{0.999}-accepting judge, but random
votes) scores \num{0.192}---identical to random play (B6).} Language contributes
only when a learned policy consumes it; a persuasive channel bolted onto an
uninformed voter is worthless. This is the paper's central result and it is
invisible to designs that do not ablate voting and language independently.

\paragraph{(2) Jointly training the social correction reverses a mismatch
artifact.} Under a train/evaluate mismatch---planner trained without the social
correction, correction switched on only at evaluation---social \emph{hurt}. Under
retrain-per-condition, where the planner is trained on $z+\delta$, the effect
flips to neutral-to-positive ($\mathrm{B1}-\mathrm{B2}=+0.020$ here; $+0.043$,
$P{=}0.90$ in a language-free replication). The integrated-social models (B0, B1)
also carry a markedly sharper world model: $\Delta z$ MSE $\approx 0.59$ versus
$3.8$--$4.8$ for the base rungs, a 6--8$\times$ reduction. The lesson is
methodological: an ablation that does not retrain per condition can invert the
sign of a component's measured effect.

\paragraph{(3) The learned planner beats random, and it scales with training.} At
medium scale the planner contrast is positive but not significant
($\mathrm{B2}-\mathrm{B6}=+0.010$). Quadrupling the offline training budget lifts
it to $+0.053$ ($P{=}0.94$) and moves the planner (B2) above the free-energy
voter (B3), indicating the factorized head is under-trained rather than
uninformative at medium scale. \note{The pending full-scale run is expected to
sharpen the planner and social contrasts; insert final numbers here.}

\subsection{Sensitivity sweeps (Table 2)}
\label{sec:sweeps}
Beyond the ablation ladder we sweep three design parameters at evaluation time
against the full trained model, holding the trained weights fixed and varying one
knob per family (Table~\ref{tab:table2}). \textbf{S1} varies the magnitude of the
social-influence correction ($\lVert\delta\rVert$ as a fraction of
$\lVert z\rVert$); $0$ disables the correction and the trained default is $0.15$.
\textbf{S2} varies the intent-bias fusion weight $\alpha$ that blends the planner's
chosen talk intent with the learned bias head when steering the decoder.
\textbf{S3} is the \emph{personality-diversity} sweep: \texttt{PERSONA\_SCALE}
controls the spread from which each agent's five-factor persona is drawn, so
$0$ yields a homogeneous population and larger values a more diverse one. Because
the persona enters both the encoder identity block and the speaker, S3 probes
whether personality \emph{diversity} per se helps or hurts coordinated villager
play---the question the paper's title raises.

\begin{table}[t]
\centering
\small
\begin{tabular}{llccccc}
\toprule
Family & Condition & Win rate \ci{95\%}{CI} & Vote acc. & Talk--vote align. & Judge acc. & $\Delta z$ MSE \\
\midrule
\multirow{3}{*}{\shortstack[l]{S1: social\\ scale}}
  & $\lVert\delta\rVert=0.00$ (off) & \TBD\ \ci{\TBD}{\TBD} & \TBD & \TBD & --    & \TBD \\
  & $\lVert\delta\rVert=0.15$       & \TBD\ \ci{\TBD}{\TBD} & \TBD & \TBD & --    & \TBD \\
  & $\lVert\delta\rVert=0.30$       & \TBD\ \ci{\TBD}{\TBD} & \TBD & \TBD & --    & \TBD \\
\midrule
\multirow{2}{*}{\shortstack[l]{S2: intent\\ fusion}}
  & $\alpha=0.3$ & \TBD\ \ci{\TBD}{\TBD} & \TBD & \TBD & \TBD & \TBD \\
  & $\alpha=0.7$ & \TBD\ \ci{\TBD}{\TBD} & \TBD & \TBD & \TBD & \TBD \\
\midrule
\multirow{3}{*}{\shortstack[l]{S3: persona\\ diversity}}
  & homogeneous ($s{=}0.0$) & \TBD\ \ci{\TBD}{\TBD} & \TBD & \TBD & \TBD & \TBD \\
  & moderate ($s{=}0.2$)    & \TBD\ \ci{\TBD}{\TBD} & \TBD & \TBD & \TBD & \TBD \\
  & diverse ($s{=}0.4$)     & \TBD\ \ci{\TBD}{\TBD} & \TBD & \TBD & \TBD & \TBD \\
\bottomrule
\end{tabular}
\caption{Table 2 --- eval-time sensitivity sweeps on the full model. \note{Populate
from \texttt{logs/sweeps/sweep\_summary.csv} after the publish run; S1 is
language-free (no judge). Report each family's trend and whether S3 shows a
diversity effect on villager win rate.}}
\label{tab:table2}
\end{table}

% ===========================================================================
\section{Discussion and Limitations}
\label{sec:discussion}
% ===========================================================================
The inseparability of language and policy (finding 1) argues against reading
strong LLM-Werewolf results as evidence of strategic competence per se: fluency
plus a judge can leave win rate at chance if the voting policy is uninformed. Our
decomposition localizes the competence in the learned planner and shows language
amplifying it rather than substituting for it.

\paragraph{Limitations.} The medium-scale contrasts for the planner and social
components are underpowered ($n=99$); the full-scale run is pending. Win rates are
governed partly by game balance (random villagers win $\approx\!0.19$ here), which
shifts with rule and correction parameters and complicates cross-paper comparison.
The judge and speaker share a model family, so judge acceptance is high
($\approx\!0.99$) and may under-penalize; a stricter or independent judge is
future work. Finally, ``truthfulness'' as scored by the judge is a proxy, not
ground-truth deception detection.

% ===========================================================================
\section{Conclusion}
% ===========================================================================
A JEPA-centered agent learns the strategic core of nine-player Werewolf, and
natural language---gated by an LLM judge---roughly doubles its win rate. But the
gain is contingent: language helps only through a learned policy that consumes it,
and an ablation must retrain per condition or risk inverting the sign of what it
measures. We release the corrected pipeline and the retrain-per-condition ladder
to support faithful component-level evaluation of language agents in social
deduction.

% ===========================================================================
\appendix
\section{Training and Implementation Details}
\label{sec:appendix}
% ===========================================================================
All components are small networks trained from self-play rollouts; the only large
models are the frozen speaker/judge (\texttt{gpt-4o-mini}). Code and configuration
are released with the paper.

\paragraph{Observation and identity block.} Each agent's observation is an
823-dimensional vector: an 808-dimensional base (public game state, per-player
features, and a MiniLM sentence embedding of recent
messages\,---\,\texttt{all-MiniLM-L6-v2}) concatenated with a 15-dimensional
\emph{identity block} (1 role bit $+$ 9 self one-hot $+$ 5 five-factor persona),
scaled by $\textsc{identity\_scale}=6.0$. The identity block conditions the latent on
who is observing and was necessary to break an encoder-collapse failure mode.

\paragraph{Networks.} The belief encoder is an MLP mapping observations to a
32-dimensional latent with an output \texttt{LayerNorm} (which bounds
$\lVert z\rVert$ and prevents collapse/explosion; a VICReg-style variance hinge,
tried first, destabilized the norm and is disabled). The world model predicts the
next latent from $(z_t, a_t)$ with an 8-dimensional one-hot action embedding
($\textsc{world\_input\_mode}=z\!\oplus\!a$). The factorized planner has three heads
(talk over five intents $\{$accuse, defend, hedge, question, vote$\}$, vote, kill)
with self/dead masking and wolves-only kill masking; planning uses top-$3$
candidates.

\paragraph{Social-influence module.} A message-conditioned network emits a
correction $\delta$ with relative scaling $\lVert\delta\rVert = s\,\lVert z\rVert$
($s=0.15$ trained), cosine trust, $\tau=0.5$, and $\lVert\delta\rVert$ bounded by a
regularizer ($\lambda_{\text{reg}}=10^{-3}$). In the \emph{integrated} regime the
planner's behavior cloning is trained on $z+\delta$; a wolf-supervision term
(weight $0.5$) trains $\delta$ to shift villager votes toward true werewolves.

\paragraph{Language and judge.} The planner emits a discrete intent and target; a
learned bias head and speaker bandit steer a frozen decoder, blended by an
intent-fusion weight $\alpha$ (default $0.65$; swept in S2). The judge
(\texttt{gpt-4o-mini}) returns JSON subscores for coherence, truthfulness,
role-consistency, and safety, and gates acceptance.

\paragraph{Training objective and hyperparameters.} Factorized training combines
behavior cloning ($\lambda_{\text{bc}}=0.5$), a talk-intent term
($\lambda_{\text{talk}}=0.5$), a small FEP-style talk-entropy regularizer
($0.01$), and the JEPA one-step latent-prediction loss, whose targets are the
re-encoded next observation under a stop-gradient. Optimizer Adam, learning rate
$10^{-3}$, batch size $64$, gradient-norm clip $1.0$. Speaker bandit and bias head
use learning rate $5\times10^{-4}$. Determinism is seeded (global seed $1337$).

\paragraph{Retrain-per-condition protocol.} The seven rungs are covered by four
deduplicated checkpoints, each written to its own namespace:
\texttt{ck\_base} (planner, no social/language; B2/B3/B5/B6),
\texttt{ck\_social} (integrated social; B1),
\texttt{ck\_llm} (language, random voting; B4), and
\texttt{ck\_full} (all components; B0). In the publication run the language-free
models (\texttt{ck\_base}, \texttt{ck\_social}) are trained at $40\times8$
games$\times$cycles for $6$ epochs, while the language models
(\texttt{ck\_llm}, \texttt{ck\_full}), which invoke the paid speaker/judge during
training, are held at $40\times2$ for $4$ epochs (the language contrast is already
significant, so the extra budget goes to the free components; games-per-cycle is
capped at $40$ to bound peak memory). The language-free rungs are evaluated at
$450$ games across seeds $\{1337, 2718, 3141\}$; the language rungs (B0, B4), whose
contrast is already significant and whose per-game cost is dominated by the paid
judge/speaker, at $150$ games over the same seeds. Sweeps (Table~2) use $300$ games
across $\{1337, 2718\}$ against \texttt{ck\_full}.

\bibliographystyle{plainnat}
\bibliography{references}

\end{document}
